\documentclass[10pt, aspectratio=169]{beamer}
% Preámbulo modular para presentaciones Beamer (Modelación Estadística)
% Diseño y paleta institucional del Tecnológico de Monterrey

\documentclass[aspectratio=169,xcolor={dvipsnames,table}]{beamer}

\usepackage[utf8]{inputenc}
\usepackage[spanish,mexico]{babel}
\usepackage{amsmath,amssymb,amsfonts,mathtools}
\usepackage{graphicx}
\usepackage{listings}
\usepackage{booktabs}
\usepackage{tikz}
\usetikzlibrary{matrix,arrows,decorations.pathmorphing,shapes.geometric,positioning}

% Paleta Institucional Tecnológico de Monterrey
\definecolor{TecAzul}{HTML}{0039A6}       % Azul TEC: color primario institucional
\definecolor{TecAzulOscuro}{HTML}{1A2E51} % Azul marino oscuro
\definecolor{TecRojo}{HTML}{EC2661}       % Rojo de acento (paleta miTec)
\definecolor{TecGris}{HTML}{646464}       % Gris neutro para comentarios y subtítulos
\definecolor{TecFondoBloque}{HTML}{F4F6F9}% Fondo suave para bloques

% Configuración de Tema Beamer
\usetheme{Boadilla}
\usecolortheme[named=TecAzulOscuro]{structure}
\setbeamercolor{alerted text}{fg=TecRojo}
\setbeamercolor{palette primary}{bg=TecAzulOscuro,fg=white}
\setbeamercolor{palette secondary}{bg=TecAzul,fg=white}
\setbeamercolor{palette tertiary}{bg=TecAzulOscuro!80!black,fg=white}
\setbeamercolor{title}{fg=TecAzulOscuro,bg=TecFondoBloque}
\setbeamercolor{frametitle}{fg=TecAzulOscuro,bg=TecFondoBloque}

% Bloques personalizados elegantes
\setbeamercolor{block title}{bg=TecAzulOscuro,fg=white}
\setbeamercolor{block body}{bg=TecFondoBloque,fg=black}
\setbeamercolor{block title alerted}{bg=TecRojo,fg=white}
\setbeamercolor{block body alerted}{bg=TecRojo!8,fg=black}
\setbeamercolor{block title example}{bg=TecAzul,fg=white}
\setbeamercolor{block body example}{bg=TecAzul!8,fg=black}

% Redondear bordes y añadir sombra sutil
\setbeamertemplate{blocks}[rounded][shadow=true]
\setbeamertemplate{navigation symbols}{}

% Pie de página limpio con número de diapositiva
\setbeamertemplate{footline}{
  \leavevmode%
  \hbox{%
  \begin{beamercolorbox}[wd=.7\paperwidth,ht=2.25ex,dp=1ex,left]{author in head/foot}%
    \hspace*{2ex}\insertshortauthor~~(\insertshortinstitute) \hfill \insertshorttitle
  \end{beamercolorbox}%
  \begin{beamercolorbox}[wd=.3\paperwidth,ht=2.25ex,dp=1ex,right]{date in head/foot}%
    \insertshortdate{}\hspace*{2em}
    \insertframenumber{} / \inserttotalframenumber\hspace*{2ex}
  \end{beamercolorbox}}%
  \vskip0pt%
}

% Configuración de Listings de Python para visualización pedagógica de código
\lstset{
    language=Python,
    basicstyle=\ttfamily\footnotesize,
    keywordstyle=\color{TecAzulOscuro}\bfseries,
    stringstyle=\color{TecRojo},
    commentstyle=\color{TecGris}\itshape,
    showstringspaces=false,
    breaklines=true,
    frame=lines,
    rulecolor=\color{TecAzulOscuro!30},
    backgroundcolor=\color{TecFondoBloque},
    aboveskip=0.5em,
    belowskip=0.5em,
    literate={á}{{\'a}}1 {é}{{\'e}}1 {í}{{\'i}}1 {ó}{{\'o}}1 {ú}{{\'u}}1
             {Á}{{\'A}}1 {É}{{\'E}}1 {Í}{{\'I}}1 {Ó}{{\'O}}1 {Ú}{{\'U}}1
             {ñ}{{\~n}}1 {Ñ}{{\~N}}1 {¿}{{?`}}1 {¡}{{!`}}1
}

% Metadatos institucionales por defecto
\title[Modelación Estadística]{Modelación Estadística para Ciencia de Datos e Ingeniería}
\author[J. Castillo Colmenares]{Juliho David Castillo Colmenares}
\institute[Tec de Monterrey]{Tecnológico de Monterrey \\ Escuela de Ingeniería y Ciencias}
\date{\today}

% Comandos y macros estadísticas compartidas para las presentaciones Beamer (Metropolis)

% Conjuntos numéricos básicos
\newcommand{\R}{\mathbb{R}}
\newcommand{\N}{\mathbb{N}}
\newcommand{\Q}{\mathbb{Q}}
\newcommand{\Z}{\mathbb{Z}}
\newcommand{\set}[1]{\left\{ #1 \right\}}
\newcommand{\abs}[1]{\left|#1\right|}
\newcommand{\norm}[1]{\left\| #1 \right\|}

% Comandos estadísticos y probabilísticos del curso
\newcommand{\Var}{\operatorname{Var}}
\newcommand{\cov}{\operatorname{Cov}}
\newcommand{\corr}{\rho}
\newcommand{\comb}[2]{\begin{pmatrix} #1 \\ #2 \end{pmatrix}}
\newcommand{\s}{\sigma}
\newcommand{\particion}{\mathfrak{P}}
\newcommand{\card}[1]{\#\left( #1 \right)}
\newcommand{\seq}[3]{\set{#1}_{#2}^{#3}}
\newcommand{\E}{\mathbb{E}}
\newcommand{\Prob}{\mathbb{P}}

% Entornos pedagógicos y cajas en español académico natural
\newenvironment{discusion}{%
  \begin{alertblock}{Pregunta de Discusión}%
}{%
  \end{alertblock}%
}

\newenvironment{reflexion}{%
  \begin{alertblock}{Reflexión para el Aula}%
}{%
  \end{alertblock}%
}

\newenvironment{paradoja}{%
  \begin{alertblock}{Paradoja de Exploración}%
}{%
  \end{alertblock}%
}

\newenvironment{motivacion}{%
  \begin{exampleblock}{Motivación: Ciencia de Datos en la Práctica}%
}{%
  \end{exampleblock}%
}

\newenvironment{retoclase}{%
  \begin{block}{Reto Rápido en Equipo}%
}{%
  \end{block}%
}


\title[Correlación]{Sección 09.01: Correlación como Premisa de la Regresión}
\subtitle{Definición, Propiedades, y la Advertencia Correlación-Causalidad}
\author[J. Castillo Colmenares]{Juliho Castillo Colmenares}
\institute{Tecnológico de Monterrey}
\date{\vspace{-1.2cm}}

\begin{document}

% Slide 1: Portada
\begin{frame}[plain]
  \vspace{-0.8cm}
  \titlepage
\end{frame}

% Slide 2: Hoja de Ruta
\begin{frame}{Hoja de Ruta --- Capítulo 09: Regresiones Lineales y Múltiples}
  \scriptsize
  ¿Dónde nos encontramos en el desarrollo modular del capítulo?
  \begin{itemize}\itemsep=0.02em
    \item \textbf{09.01} Correlación como Premisa de la Regresión \emph{(Hoy)}
    \item \textbf{09.02} Introducción a la Regresión Lineal
    \item \textbf{09.03} Matemáticas de la Regresión: Derivación MCO y $R^2$
    \item \textbf{09.04} Regresión Lineal sobre Datos Simulados
    \item \textbf{09.05} Coeficientes Óptimos, Pruebas $t$/$F$ y RSE
    \item \textbf{09.06} Implementación en Python con \texttt{statsmodels}
    \item \textbf{09.07} Regresión Múltiple, Ecuación Normal y Ridge/Lasso
    \item \textbf{09.08} Validación de Modelos y $k$-fold Cross-Validation
    \item \textbf{09.09} Diagnóstico de Residuos y Multicolinealidad (VIF)
    \item \textbf{09.10} Regresión con \texttt{scikit-learn}
    \item \textbf{09.11} Variables Categóricas y Variables Muda
    \item \textbf{09.12} Transformaciones No Lineales y Regresión Polinomial
  \end{itemize}
  \pause
  \vspace{-0.05cm}
  \begin{block}{Objetivo de la Sesión}
    Establecer la correlación como la premisa conceptual de toda la modelación predictiva: definir formalmente el coeficiente de Pearson, estudiar sus propiedades y rango, y advertir rigurosamente que correlación \textbf{no} implica causalidad --- la distinción que abrirá la puerta a la regresión lineal en la Sección 09.02.
  \end{block}
\end{frame}

% Slide 3: Motivación
\begin{frame}{¿Por Qué Empezar con Correlación?}
  \footnotesize
  \begin{motivacion}
    La correlación es la premisa de la modelación predictiva: es el factor con base en el cual podemos anticipar el comportamiento de una variable a partir de otra. Si $X$ y $Y$ están fuertemente correlacionadas, un cambio en una sugiere un cambio asociado en la otra.
  \end{motivacion}

  \vspace{0.15cm}
  \pause
  \begin{itemize}\itemsep=0.1cm
    \item Una buena correlación asegura una relación matemática funcional aproximada $Y=f(X)$. \pause
    \item Esa relación puede ser de cualquier forma: lineal ($y=ax+b$), exponencial ($y=be^{ax}$), u otra --- la regresión lineal se enfoca en el caso lineal.
  \end{itemize}
\end{frame}

% Slide 4: Definición formal
\begin{frame}{Definición Formal del Coeficiente de Correlación}
  \footnotesize
  \begin{block}{Coeficiente de Correlación de Pearson}
    Para variables $X,Y$ con promedios muestrales $\bar x,\bar y$:
    $$r_{XY}=\frac{\sum_i(x_i-\bar x)(y_i-\bar y)}{\sqrt{\sum_i(x_i-\bar x)^2\sum_i(y_i-\bar y)^2}}=\frac{S_{XY}}{\sqrt{S_{XX}S_{YY}}}$$
  \end{block}
  \pause

  \vspace{0.1cm}
  \begin{alertblock}{Interpretación}
    $r_{XY}$ normaliza la covarianza $S_{XY}$ por las dispersiones individuales de $X$ y $Y$, produciendo una medida adimensional de asociación lineal comparable entre cualquier par de variables.
  \end{alertblock}
\end{frame}

% Slide 5: Propiedades
\begin{frame}{Propiedades del Coeficiente de Correlación}
  \footnotesize
  \begin{itemize}\itemsep=0.12cm
    \item \textbf{Rango acotado:} $-1\leq r_{XY}\leq1$, consecuencia directa de la desigualdad de Cauchy-Schwarz. \pause
    \item \textbf{Signo:} $r_{XY}>0$ indica relación directa (ambas variables suben o bajan juntas); $r_{XY}<0$ indica relación inversa. \pause
    \item \textbf{Magnitud:} entre mayor $|r_{XY}|$, más fuerte y predecible es la relación lineal; $|r_{XY}|=1$ indica una relación lineal exacta.
  \end{itemize}
\end{frame}

% Slide 6: Advertencia
\begin{frame}{¡Advertencia! Correlación No Implica Causalidad}
  \footnotesize
  \begin{alertblock}{Observación Crítica}
    Una correlación fuerte sugiere una relación útil para la predicción, pero \textbf{no implica que esa relación sea el único factor} que explica el comportamiento observado, ni que una variable \emph{cause} a la otra.
  \end{alertblock}
  \pause

  \vspace{0.1cm}
  \begin{block}{Ejemplo Clásico: la Variable de Confusión (\emph{Confounder})}
    Existe una correlación positiva muy fuerte entre el número de palabras que conoce un niño y el número de calzado que usa. Ninguna causa a la otra: ambas son consecuencia de una tercera variable oculta, la \textbf{edad}.
  \end{block}
\end{frame}

% Slide 7: El caso publicitario
\begin{frame}{Caso de Estudio: Publicidad y Ventas}
  \footnotesize
  Conjunto de datos clásico: gasto en publicidad por medio (\texttt{TV}, \texttt{Radio}, \texttt{Newspaper}) vs. \texttt{Sales} de un producto.
  \pause

  \vspace{0.1cm}
  \begin{block}{Coeficientes de Correlación Observados}
    $$r(\texttt{TV},\texttt{Sales})\approx0.782, \qquad r(\texttt{Radio},\texttt{Sales})\approx0.576$$
  \end{block}
  \pause

  \vspace{0.1cm}
  \begin{alertblock}{Lectura Preliminar}
    \texttt{TV} exhibe la asociación lineal más fuerte con \texttt{Sales}: un candidato natural para el primer predictor de un modelo de regresión (Secciones 09.02 en adelante).
  \end{alertblock}
\end{frame}

% Slide 8: Lab Python (Bloque 1)
\begin{frame}[fragile]{Laboratorio en Python: Coeficiente de Pearson desde su Definición}
  \scriptsize
  Cálculo directo de $r_{XY}$ a partir de sumas de desviaciones, verificado contra \texttt{np.corrcoef} (\texttt{09.01\_correlation.py}):
  {\lstset{basicstyle=\fontsize{3.6pt}{4.3pt}\selectfont\ttfamily,aboveskip=0.05em,belowskip=0.05em}
  \lstinputlisting[firstline=17, lastline=27]{../../code/09_regresiones/09.01_correlation.py}}
\end{frame}

% Slide 9: Lab Python (Bloque 2)
\begin{frame}[fragile]{Laboratorio en Python: Rango y Fuerza de la Correlación}
  \scriptsize
  Seis escenarios ilustrando el rango $[-1,1]$ y la interpretación de signo/magnitud:
  {\lstset{basicstyle=\fontsize{3.4pt}{4.1pt}\selectfont\ttfamily,aboveskip=0.05em,belowskip=0.05em}
  \lstinputlisting[firstline=30, lastline=45]{../../code/09_regresiones/09.01_correlation.py}}
\end{frame}

% Slide 10: Lab Python (Bloque 3)
\begin{frame}[fragile]{Laboratorio en Python: Correlación No Es Causalidad}
  \scriptsize
  Correlación espuria entre tamaño de calzado y vocabulario, resuelta vía correlación parcial controlando por edad:
  {\lstset{basicstyle=\fontsize{3.2pt}{3.9pt}\selectfont\ttfamily,aboveskip=0.05em,belowskip=0.05em}
  \lstinputlisting[firstline=48, lastline=64]{../../code/09_regresiones/09.01_correlation.py}}
\end{frame}

% Slide 11: Salida Terminal
\begin{frame}[fragile]{Laboratorio en Python: Salida en Terminal}
  \scriptsize
  Salida estándar generada al ejecutar el script del laboratorio en Python:
  \vspace{0.02cm}
  \begin{block}{Salida Estándar en Consola}
    \fontsize{3.8pt}{4.6pt}\selectfont
    \begin{verbatim}
=== Block 1: Pearson from Definition ===
r (from definition) = 0.8959 (matches np.corrcoef)

=== Block 2: Range and Strength ===
perfect positive: r=+1.0000   strong positive: r=+0.9638
weak positive:     r=+0.3710   no relationship: r=-0.1333
strong negative:   r=-0.9690   perfect negative: r=-1.0000

=== Block 3: Correlation is Not Causation ===
Raw correlation(shoe size, vocabulary) = 0.9517 (spurious)
Partial correlation(... | age)         = 0.0486 (collapses)
    \end{verbatim}
  \end{block}
\end{frame}

% Slide 12: Discusión en clase (Enunciado)
\begin{frame}{Discusión en Clase --- El Peligro de un Solo Coeficiente}
  \footnotesize
  \begin{block}{Pregunta}
    En el caso publicitario, $r(\texttt{Radio},\texttt{Sales})\approx0.576$ es una correlación moderada, no débil. ¿Por qué no basta con mirar $r$ para decidir si \texttt{Radio} debe incluirse en un modelo de regresión múltiple junto con \texttt{TV}?
  \end{block}

  \vspace{0.15cm}
  \pause
  \begin{alertblock}{Pista}
    Piensa en si \texttt{Radio} y \texttt{TV} podrían estar, a su vez, correlacionadas entre sí.
  \end{alertblock}
\end{frame}

% Slide 13: Discusión en clase (Respuesta)
\begin{frame}{Discusión en Clase --- Respuesta}
  \scriptsize
  La correlación bivariada $r(\texttt{Radio},\texttt{Sales})$ ignora por completo la posible relación entre \texttt{Radio} y \texttt{TV}: \pause

  \vspace{0.1cm}
  \begin{alertblock}{Interpretación}
    Si \texttt{Radio} y \texttt{TV} comparten variabilidad (por ejemplo, campañas que combinan ambos medios simultáneamente), parte de lo que $r(\texttt{Radio},\texttt{Sales})$ captura ya está siendo explicado por \texttt{TV}. Esta es exactamente la razón por la que la \textbf{regresión múltiple} (Sección 09.07) --- y no una colección de correlaciones bivariadas aisladas --- es la herramienta correcta para cuantificar el efecto de cada predictor \emph{controlando por los demás}.
  \end{alertblock}
\end{frame}

% Slide 14: Cierre
\begin{frame}{Cierre de Sección: Correlación}
  \begin{reflexion}
    La correlación cuantifica la fuerza y dirección de una asociación lineal, y es la motivación conceptual de todo el capítulo: si dos variables no están relacionadas, ningún modelo de regresión podrá predecir la una a partir de la otra con utilidad práctica. Pero la correlación describe asociación, no mecanismo causal --- una distinción que todo científico de datos debe mantener presente.
  \end{reflexion}

  \vspace{0.2cm}
  \pause
  \begin{block}{Lecciones Clave}
    \begin{itemize}\itemsep=0.08cm
      \item \textbf{Rango:} $r_{XY}\in[-1,1]$, con signo indicando dirección y magnitud indicando fuerza.
      \item \textbf{Correlación $\neq$ causalidad:} una tercera variable de confusión puede generar asociación espuria.
      \item \textbf{Correlación bivariada $\neq$ efecto controlado:} predictores correlacionados entre sí requieren regresión múltiple, no correlaciones aisladas.
    \end{itemize}
  \end{block}
\end{frame}

% Slide 15: Perspectiva Modular
\begin{frame}{Perspectiva Modular --- El Próximo Reto}
  \scriptsize
  \begin{retoclase}
    La correlación mide asociación pero no permite predecir un valor numérico específico de $Y$ dado un valor de $X$. La Sección 09.02 introduce la \textbf{regresión lineal}: el modelo $Y_e=\hat\beta_0+\hat\beta_1X$ que convierte la asociación medida por $r_{XY}$ en una herramienta predictiva concreta.
  \end{retoclase}

  \vspace{0.08cm}
  \pause
  \begin{alertblock}{Preparación Recomendada}
    \begin{itemize}\itemsep=0.06cm
      \item Reflexiona sobre qué información adicional (más allá de $r_{XY}$) necesitarías para predecir un valor concreto de \texttt{Sales} dado un nivel de gasto en \texttt{TV}.
      \item Repasa la ecuación de una recta $y=ax+b$: la regresión lineal estimará $a$ y $b$ a partir de los datos.
    \end{itemize}
  \end{alertblock}
\end{frame}

\end{document}
